% GENERATED FILE. Edit canonical lessons, then run npm run generate:books.
% canonical-source-hash: fnv1a64:324d8b10a91ad7ca
% canonical-lessons: FR-C333-levure, FR-C333-petrir, FR-C333-rouleau-a-patisserie, FR-C333-etaler, FR-C333-demouler, FR-R333-revision-332-333, FR-R333-revision-final-333

\chapter{Faire du pain}
\label{ch:fr-faire-du-pain}
\hlchaptermodality{333}

\begin{chapteropening}
\textbf{By the end of this chapter:} \emph{I can describe making dough: yeast, kneading, the rolling pin, rolling out and turning out a cake.}

Complete the last lesson of chapter 333: I can describe making dough: yeast, kneading, the rolling pin, rolling out and turning out a cake.
\end{chapteropening}

\section[la levure]{la levure — the yeast; baking powder}

\label{lesson:FR-C333-levure}



\begin{quote}
Before the new one: say the French for to season, then the French for to taste.
\end{quote}

\subsection*{You'll want to know: la levure}
\textbf{la levure} — \textquotedblleft{}the yeast; baking powder\textquotedblright{}.

\begin{quote}
Le boulanger met de la levure dans la pâte.
\end{quote}

\begin{quote}
\emph{The baker puts yeast in the dough.}
\end{quote}

\emph{La levure chimique} — baking powder.

\begin{cousinweb}
\emph{levure} is built on \emph{lever} \textquotedblleft{}to raise\textquotedblright{}, which is inherited from Latin \emph{levāre} \textquotedblleft{}to lift\textquotedblright{} (EtymDB).
\end{cousinweb}

\begin{grammarlens}[title={What you've built}]
One more word for telling, explaining and describing at B1.
\end{grammarlens}

\subsection*{Your turn}
\begin{itemize}
\raggedright
  \item \emph{Say it:} \emph{la levure}
  \item \emph{Say it:} \emph{la levure}, once more
  \item \emph{Say it:} a sentence of your own with \emph{la levure}
  \item \emph{Recall:} say the French for to season, then the French for to taste, then say \emph{la levure} again
\end{itemize}

\begin{tcolorbox}[breakable,colback=teal!4,colframe=teal!35!black,title={Before you move on}]
What does la levure mean? (\textquotedblleft{}The yeast; baking powder\textquotedblright{}.) Say it once more.
\end{tcolorbox}
\section[pétrir]{pétrir — to knead}

\label{lesson:FR-C333-petrir}



\begin{quote}
Before the new one: say the French for to taste, then the French for the yeast.
\end{quote}

\subsection*{You'll want to know: pétrir}
\textbf{pétrir} — \textquotedblleft{}to knead\textquotedblright{}.

\begin{quote}
Le boulanger pétrit la pâte pendant dix minutes.
\end{quote}

\begin{quote}
\emph{The baker kneads the dough for ten minutes.}
\end{quote}

Regular, like \emph{finir}: \emph{il pétrit}, \emph{nous pétrissons}.

\begin{cousinweb}
\emph{pétrir} is usually traced to Late Latin \emph{pistrīre} \textquotedblleft{}to knead\textquotedblright{}, built on \emph{pistor} \textquotedblleft{}a miller, a baker\textquotedblright{} (Lewis \& Short); EtymDB lists no parent.
\end{cousinweb}

\begin{grammarlens}[title={What you've built}]
One more word for telling, explaining and describing at B1.
\end{grammarlens}

\subsection*{Your turn}
\begin{itemize}
\raggedright
  \item \emph{Say it:} \emph{pétrir}
  \item \emph{Say it:} \emph{pétrir}, once more
  \item \emph{Say it:} a sentence of your own with \emph{pétrir}
  \item \emph{Recall:} say the French for to taste, then the French for the yeast, then say \emph{pétrir} again
\end{itemize}

\begin{tcolorbox}[breakable,colback=teal!4,colframe=teal!35!black,title={Before you move on}]
What does pétrir mean? (\textquotedblleft{}To knead\textquotedblright{}.) Say it once more.
\end{tcolorbox}
\section[le rouleau à pâtisserie]{le rouleau à pâtisserie — the rolling pin}

\label{lesson:FR-C333-rouleau-a-patisserie}



\begin{quote}
Before the new one: say the French for the yeast, then the French for to knead.
\end{quote}

\subsection*{You'll want to know: le rouleau à pâtisserie}
\textbf{le rouleau à pâtisserie} — \textquotedblleft{}the rolling pin\textquotedblright{}.

\begin{quote}
Le rouleau à pâtisserie est dans le tiroir.
\end{quote}

\begin{quote}
\emph{The rolling pin is in the drawer.}
\end{quote}

\emph{Un rouleau} alone is a roll: \emph{un rouleau de papier}.

\begin{cousinweb}
\emph{rouleau} is built on \emph{rôle} \textquotedblleft{}a roll\textquotedblright{}, from Latin \emph{rotulus} \textquotedblleft{}a little wheel\textquotedblright{}, built on \emph{rota} (Lewis \& Short); \emph{pâtisserie} is built on \emph{pâte}.
\end{cousinweb}

\begin{grammarlens}[title={What you've built}]
One more word for telling, explaining and describing at B1.
\end{grammarlens}

\subsection*{Your turn}
\begin{itemize}
\raggedright
  \item \emph{Say it:} \emph{le rouleau à pâtisserie}
  \item \emph{Say it:} \emph{le rouleau à pâtisserie}, once more
  \item \emph{Say it:} a sentence of your own with \emph{le rouleau à pâtisserie}
  \item \emph{Recall:} say the French for the yeast, then the French for to knead, then say \emph{le rouleau à pâtisserie} again
\end{itemize}

\begin{tcolorbox}[breakable,colback=teal!4,colframe=teal!35!black,title={Before you move on}]
What does le rouleau à pâtisserie mean? (\textquotedblleft{}The rolling pin\textquotedblright{}.) Say it once more.
\end{tcolorbox}
\section[étaler]{étaler — to roll out, to spread}

\label{lesson:FR-C333-etaler}



\begin{quote}
Before the new one: say the French for to knead, then the French for the rolling pin.
\end{quote}

\subsection*{You'll want to know: étaler}
\textbf{étaler} — \textquotedblleft{}to roll out, to spread\textquotedblright{}.

\begin{quote}
Il étale la pâte avec le rouleau à pâtisserie.
\end{quote}

\begin{quote}
\emph{He rolls out the dough with the rolling pin.}
\end{quote}

Also to spread butter: \emph{étaler du beurre sur le pain}.

\begin{cousinweb}
\emph{étaler} is built on Old French \emph{estal} \textquotedblleft{}a place, a stall\textquotedblright{}, which is usually traced to a Frankish word (EtymDB lists no parent): to lay out on a stall.
\end{cousinweb}

\begin{grammarlens}[title={What you've built}]
One more word for telling, explaining and describing at B1.
\end{grammarlens}

\subsection*{Your turn}
\begin{itemize}
\raggedright
  \item \emph{Say it:} \emph{étaler}
  \item \emph{Say it:} \emph{étaler}, once more
  \item \emph{Say it:} a sentence of your own with \emph{étaler}
  \item \emph{Recall:} say the French for to knead, then the French for the rolling pin, then say \emph{étaler} again
\end{itemize}

\begin{tcolorbox}[breakable,colback=teal!4,colframe=teal!35!black,title={Before you move on}]
What does étaler mean? (\textquotedblleft{}To roll out, to spread\textquotedblright{}.) Say it once more.
\end{tcolorbox}
\section[démouler]{démouler — to turn out (of a tin)}

\label{lesson:FR-C333-demouler}



\begin{quote}
Before the new one: say the French for the rolling pin, then the French for to roll out.
\end{quote}

\subsection*{You'll want to know: démouler}
\textbf{démouler} — \textquotedblleft{}to turn out (of a tin)\textquotedblright{}.

\begin{quote}
Elle démoule le gâteau quand il est froid.
\end{quote}

\begin{quote}
\emph{She turns out the cake when it is cold.}
\end{quote}

The opposite of putting into \emph{le moule}.

\begin{cousinweb}
\emph{démouler} is built on \emph{dé-} and \emph{moule} \textquotedblleft{}mould\textquotedblright{}, which comes from Latin \emph{modulus} \textquotedblleft{}a small measure\textquotedblright{} (EtymDB).
\end{cousinweb}

\begin{grammarlens}[title={What you've built}]
One more word for telling, explaining and describing at B1.
\end{grammarlens}

\subsection*{Your turn}
\begin{itemize}
\raggedright
  \item \emph{Say it:} \emph{démouler}
  \item \emph{Say it:} \emph{démouler}, once more
  \item \emph{Say it:} a sentence of your own with \emph{démouler}
  \item \emph{Recall:} say the French for the rolling pin, then the French for to roll out, then say \emph{démouler} again
\end{itemize}

\begin{tcolorbox}[breakable,colback=teal!4,colframe=teal!35!black,title={Before you move on}]
What does démouler mean? (\textquotedblleft{}To turn out (of a tin)\textquotedblright{}.) Say it once more.
\end{tcolorbox}
\section[(dialogue)]{Pas à pas, encore une fois}

\label{lesson:FR-R333-revision-332-333}



\begin{quote}
Say the French for the colander and to drain.
\end{quote}

\subsection*{Your turn}
Say these aloud:

\begin{itemize}
\raggedright
  \item the colander, to drain, to sprinkle, to season, to taste
  \item the yeast, to knead, the rolling pin, to roll out, to turn out
\end{itemize}

\begin{tcolorbox}[breakable,colback=teal!4,colframe=teal!35!black,title={Before you move on}]
The colander? (\textbf{la passoire}.) To drain? (\textbf{égoutter}.) To sprinkle? (\textbf{saupoudrer}.) To season? (\textbf{assaisonner}.) To taste? (\textbf{déguster}.) The yeast? (\textbf{la levure}.) To knead? (\textbf{pétrir}.) The rolling pin? (\textbf{le rouleau à pâtisserie}.) To roll out? (\textbf{étaler}.) To turn out? (\textbf{démouler}.)
\end{tcolorbox}
\section[(dialogue)]{Faire du pain — once more}

\label{lesson:FR-R333-revision-final-333}



\begin{quote}
Say the French for the yeast and to knead.
\end{quote}

\subsection*{Your turn}
\begin{itemize}
\raggedright
  \item \emph{Say it:} the yeast, to knead, the rolling pin, to roll out, to turn out
\end{itemize}

\begin{tcolorbox}[breakable,colback=teal!4,colframe=teal!35!black,title={Before you move on}]
The yeast? (\textbf{la levure}.) To knead? (\textbf{pétrir}.) The rolling pin? (\textbf{le rouleau à pâtisserie}.) To roll out? (\textbf{étaler}.) To turn out? (\textbf{démouler}.)
\end{tcolorbox}
